\chapter{Действительные (вещественные) числа}

% -----------------------------------------------------------------------------
\lecture[Аксиоматика и свойства действительных чисел]{Аксиоматика и некоторые общие свойства множества действительных чисел}

\subsection{Определение множества действительных чисел}

\begin{definition}{Множество действительных чисел}{real-numbers}
  Множество $\R$ называется множеством действительных (вещественных) чисел,
  а его элементы — действительными (вещественными) числами, если выполнен
  следующий комплекс условий, называемый аксиоматикой вещественных чисел.
\end{definition}

\begin{axiom}{Сложения}{addition-axioms}
  Определено отображение (операция сложения)
  \[
    +\colon \R\times\R\to\R,
  \]
  сопоставляющее каждой упорядоченной паре $(x,y)$ элементов $x,y$ из $\R$
  некоторый элемент $x+y\in\R$, называемый суммой $x$ и $y$. При этом
  выполнены следующие условия:
  \begin{enumerate}[label=\arabic*${}^{+}$.]
    \item Существует нейтральный элемент $0$ (называемый в случае сложения
      нулем) такой, что для любого $x\in\R$
      \[
        x+0=0+x=x.
      \]\par
    \item Для любого элемента $x\in\R$ имеется элемент $-x\in\R$,
      называемый противоположным к $x$, такой, что
      \[
        x+(-x)=(-x)+x=0.
      \]\par
    \item Операция $+$ ассоциативна, т.~е. для любых элементов $x,y,z$ из
      $\R$ выполнено
      \[
        x+(y+z)=(x+y)+z.
      \]\par
    \item Операция $+$ коммутативна, т.~е. для любых элементов $x,y$ из
      $\R$ выполнено
      \[
        x+y=y+x.
      \]
  \end{enumerate}

  Если на каком-то множестве $G$ определена операция, удовлетворяющая
  аксиомам $1^{+}$, $2^{+}$, $3^{+}$, то говорят, что на $G$ задана структура
  группы или что $G$ есть группа. Если операцию называют сложением, то
  группу называют аддитивной. Если, кроме того, известно, что операция
  коммутативна, т.~е. выполнено условие $4^{+}$, то группу называют
  коммутативной или абелевой.
\end{axiom}

\begin{axiom}{Умножения}{multiplication-axioms}
  Определено отображение (операция умножения)
  \[
    \cdot\colon \R\times\R\to\R,
  \]
  сопоставляющее каждой упорядоченной паре $(x,y)$ элементов $x,y$ из $\R$
  некоторый элемент $x\cdot y\in\R$, называемый произведением $x$ и $y$,
  причем так, что выполнены следующие условия:
  \begin{enumerate}[label=\arabic*${}^{\cdot}$.]
    \item Существует нейтральный элемент $1\in\R\setminus\{0\}$
      (называемый в случае умножения единицей) такой, что $\forall x\in\R$
      \[
        x\cdot1=1\cdot x=x.
      \]\par
    \item Для любого элемента $x\in\R\setminus\{0\}$ имеется элемент
      $x^{-1}\in\R$, называемый обратным, такой, что
      \[
        x\cdot x^{-1}=x^{-1}\cdot x=1.
      \]\par
    \item Операция $\cdot$ ассоциативна, т.~е. для любых $x,y,z$ из $\R$
      \[
        x\cdot(y\cdot z)=(x\cdot y)\cdot z.
      \]\par
    \item Операция $\cdot$ коммутативна, т.~е. для любых $x,y$ из $\R$
      \[
        x\cdot y=y\cdot x.
      \]
  \end{enumerate}
\end{axiom}

\begin{axiom}{Связь сложения и умножения}{addition-multiplication-axiom}
  Умножение дистрибутивно по отношению к сложению, т.~е.
  $\forall x,y,z\in\R$
  \[
    (x+y)z=xz+yz.
  \]
\end{axiom}

Если на каком-то множестве $G$ действуют две операции, удовлетворяющие
всем перечисленным аксиомам, то $G$ называется алгебраическим полем или
просто полем.

\begin{axiom}{Порядка}{order-axioms}
  Между элементами $\R$ имеется отношение
  $\le$, т.~е. для элементов $x,y$ из $\R$ установлено, выполняется ли
  $x\le y$ или нет. При этом должны удовлетворяться следующие условия:
  \begin{enumerate}[leftmargin=3em]
    \item[$0_{\le}$.] $\forall x\in\R\;(x\le x)$.\par
    \item[$1_{\le}$.] $(x\le y)\land(y\le x)\Rightarrow(x=y)$.\par
    \item[$2_{\le}$.] $(x\le y)\land(y\le z)\Rightarrow(x\le z)$.\par
    \item[$3_{\le}$.] $\forall x\in\R\;\forall y\in\R\;(x\le y)\lor(y\le x)$.
  \end{enumerate}

  Отношение $\le$ в $\R$ называется отношением неравенства. Множество,
  между некоторыми элементами которого имеется отношение, удовлетворяющее
  аксиомам $0_{\le}$, $1_{\le}$, $2_{\le}$, называют частично упорядоченным,
  а если, сверх того, выполнена аксиома $3_{\le}$, т.~е. любые два элемента
  множества сравнимы, то множество называется линейно упорядоченным.
\end{axiom}

\begin{axiom}{Связь сложения и порядка в $\R$}{addition-order-axiom}
  Если $x,y,z$ — элементы $\R$, то
  \[
    (x\le y)\Rightarrow(x+z\le y+z).
  \]
\end{axiom}

\begin{axiom}{Связь умножения и порядка в $\R$}{multiplication-order-axiom}
  Если $x,y$ — элементы $\R$, то
  \[
    (0\le x)\land(0\le y)\Rightarrow(0\le x\cdot y).
  \]
\end{axiom}

\begin{axiom}{Полноты (непрерывности)}{completeness-axiom}
  Если $X$ и $Y$ — непустые
  подмножества $\R$, обладающие тем свойством, что для любых элементов
  $x\in X$ и $y\in Y$ выполнено $x\le y$, то существует такое $c\in\R$,
  что $x\le c\le y$ для любых элементов $x\in X$ и $y\in Y$.
\end{axiom}

\subsection{Некоторые общие алгебраические свойства действительных чисел}

\begin{properties}{Следствия аксиом сложения}{addition-consequences}
  \begin{enumerate}[label=\arabic*${}^{\circ}$.]
    \item В множестве действительных чисел имеется только один нуль.\par
    \item В множестве действительных чисел у каждого элемента имеется
      единственный противоположный элемент.\par
    \item Уравнение
      \[
        a+x=b
      \]
      в $\R$ имеет и притом единственное решение
      \[
        x=b+(-a).
      \]
  \end{enumerate}
\end{properties}

\begin{properties}{Следствия аксиом умножения}{multiplication-consequences}
  \begin{enumerate}[label=\arabic*${}^{\circ}$.]
    \item В множестве действительных чисел имеется только одна единица.\par
    \item Для каждого числа $x\ne0$ имеется только один обратный элемент
      $x^{-1}$.\par
    \item Уравнение $a\cdot x=b$ при $a\in\R\setminus\{0\}$ имеет и притом
      единственное решение $x=b\cdot a^{-1}$.
  \end{enumerate}
\end{properties}

\begin{properties}{Следствия аксиомы связи сложения и умножения}{field-consequences}
  \begin{enumerate}[label=\arabic*${}^{\circ}$.]
    \item Для любого $x\in\R$
      \[
        x\cdot0=0\cdot x=0.
      \]\par
    \item
      \[
        (x\cdot y=0)\Rightarrow(x=0)\lor(y=0).
      \]\par
    \item Для любого $x\in\R$
      \[
        -x=(-1)\cdot x.
      \]\par
    \item Для любого числа $x\in\R$
      \[
        (-1)(-x)=x.
      \]\par
    \item Для любого числа $x\in\R$
      \[
        (-x)(-x)=x\cdot x.
      \]
  \end{enumerate}
\end{properties}

\begin{definition}{Строгое неравенство}{strict-inequality}
  Отношение $x\le y$ записывают также в виде $y\ge x$; отношение $x\le y$
  при $x\ne y$ записывают в виде $x<y$ или в виде $y>x$ и называют строгим
  неравенством.
\end{definition}

\begin{properties}{Следствия аксиом порядка}{order-consequences}
  \begin{enumerate}[label=\arabic*${}^{\circ}$.]
    \item Для любых $x,y\in\R$ всегда имеет место в точности одно из
      соотношений:
      \[
        x<y,\qquad x=y,\qquad x>y.
      \]\par
    \item Для любых чисел $x,y,z$ из $\R$
      \[
        (x<y)\land(y\le z)\Rightarrow(x<z),
      \]
      \[
        (x\le y)\land(y<z)\Rightarrow(x<z).
      \]
  \end{enumerate}
\end{properties}

\begin{properties}{Следствия аксиом связи порядка со сложением и умножением}{ordered-field-consequences}
  \begin{enumerate}[label=\arabic*${}^{\circ}$.]
    \item Для любых чисел $x,y,z,w$ из $\R$
      \[
        \begin{aligned}
          (x<y)&\Rightarrow(x+z<y+z),\\
          (0<x)&\Rightarrow(-x<0),\\
          (x\le y)\land(z\le w)&\Rightarrow(x+z\le y+w),\\
          (x\le y)\land(z<w)&\Rightarrow(x+z<y+w).
        \end{aligned}
      \]\par
    \item Если $x,y,z$ — числа из $\R$, то
      \[
        \begin{aligned}
          (0<x)\land(0<y)&\Rightarrow(0<xy),\\
          (x<0)\land(y<0)&\Rightarrow(0<xy),\\
          (x<0)\land(0<y)&\Rightarrow(xy<0),\\
          (x<y)\land(0<z)&\Rightarrow(xz<yz),\\
          (x<y)\land(z<0)&\Rightarrow(yz<xz).
        \end{aligned}
      \]\par
    \item $0<1$.\par
    \item
      \[
        (0<x)\Rightarrow(0<x^{-1})
      \]
      и
      \[
        (0<x)\land(x<y)\Rightarrow(0<y^{-1})\land(y^{-1}<x^{-1}).
      \]
  \end{enumerate}
\end{properties}

\subsection{Аксиома полноты и существование верхней (нижней) грани числового множества}

\begin{definition}{Ограниченное сверху и снизу множество}{bounded-one-sided}
  Говорят, что множество $X\subset\R$ ограничено сверху (снизу), если
  существует число $c\in\R$ такое, что $x\le c$ (соответственно, $c\le x$)
  для любого $x\in X$.

  Число $c$ в этом случае называют верхней (соответственно, нижней) границей
  множества $X$ или также мажорантой (минорантой) множества $X$.
\end{definition}

\begin{definition}{Ограниченное множество}{bounded-set}
  Множество, ограниченное и сверху и снизу, называется ограниченным.
\end{definition}

\begin{definition}{Максимальный и минимальный элементы}{max-min-elements}
  Элемент $a\in X$ называется наибольшим или максимальным (наименьшим или
  минимальным) элементом множества $X\subset\R$, если $x\le a$
  (соответственно, $a\le x$) для любого элемента $x\in X$.
  \[
    (a=\max X):=(a\in X\land\forall x\in X\;(x\le a)),
  \]
  \[
    (a=\min X):=(a\in X\land\forall x\in X\;(a\le x)).
  \]
\end{definition}

\begin{properties}{Единственность максимального и минимального элементов}{max-min-uniqueness}
  Если в числовом множестве есть максимальный (минимальный) элемент, то он
  только один.
\end{properties}

\begin{definition}{Точная верхняя граница}{supremum}
  Наименьшее из чисел, ограничивающих множество $X\subset\R$ сверху,
  называется верхней гранью (или точной верхней границей) множества $X$ и
  обозначается $\sup X$.
  \[
    (s=\sup X):=
    \forall x\in X\bigl((x\le s)\land
      (\forall s'<s\;\exists x'\in X\;(s'<x'))\bigr).
  \]
  Таким образом,
  \[
    \sup X:=\min\{c\in\R\mid\forall x\in X\;(x\le c)\}.
  \]
\end{definition}

\begin{definition}{Точная нижняя граница}{infimum}
  Нижняя грань (точная нижняя граница) множества $X$ есть наибольшая из
  нижних границ множества $X$.
  \[
    (i=\inf X):=
    \forall x\in X\bigl((i\le x)\land
      (\forall i'>i\;\exists x'\in X\;(x'<i'))\bigr).
  \]
  Таким образом,
  \[
    \inf X:=\max\{c\in\R\mid\forall x\in X\;(c\le x)\}.
  \]
\end{definition}

\begin{lemma}{Принцип верхней грани}{upper-bound-principle}
  Всякое непустое ограниченное сверху подмножество множества вещественных
  чисел имеет и притом единственную верхнюю грань.
\end{lemma}

\begin{lemma}{Принцип нижней грани}{lower-bound-principle}
  \[
    (X\text{ не пусто и ограничено снизу})\Rightarrow(\exists!\inf X).
  \]
\end{lemma}

% -----------------------------------------------------------------------------
\lecture[Важнейшие классы действительных чисел]{Важнейшие классы действительных чисел и вычислительные аспекты операций с действительными числами}

\subsection{Натуральные числа и принцип математической индукции}

\begin{definition}{Индуктивное множество}{inductive-set-real}
  Множество $X\subset\R$ называется индуктивным, если вместе с каждым
  числом $x\in X$ ему принадлежит также число $x+1$.
\end{definition}

\begin{properties}{Пересечение индуктивных множеств}{intersection-inductive}
  Пересечение
  \[
    X=\bigcap_{\alpha\in A}X_{\alpha}
  \]
  любого семейства индуктивных множеств $X_{\alpha}$, если оно непусто,
  является индуктивным множеством.
\end{properties}

\begin{definition}{Множество натуральных чисел}{natural-numbers}
  Множеством натуральных чисел называется наименьшее индуктивное множество,
  содержащее $1$, т.~е. пересечение всех индуктивных множеств, содержащих
  число $1$.

  Множество натуральных чисел обозначают символом $\N$; его элементы
  называются натуральными числами.
\end{definition}

\begin{properties}{Принцип математической индукции}{induction-principle}
  Если подмножество $E$ множества натуральных чисел $\N$ таково, что $1\in E$
  и вместе с числом $x\in E$ множеству $E$ принадлежит число $x+1$, то
  $E=\N$.
  \[
    (E\subset\N)\land(1\in E)\land
    \bigl(\forall x\in E\;(x\in E\Rightarrow(x+1)\in E)\bigr)
    \Rightarrow E=\N.
  \]
\end{properties}

\begin{properties}{Свойства натуральных чисел}{natural-properties}
  \begin{enumerate}[label=\arabic*${}^{\circ}$.]
    \item Сумма и произведение натуральных чисел являются натуральными
      числами.\par
    \item
      \[
        (n\in\N)\land(n\ne1)\Rightarrow((n-1)\in\N).
      \]\par
    \item Для любого $n\in\N$ в множестве $\{x\in\N\mid n<x\}$ есть
      минимальный элемент, причем
      \[
        \min\{x\in\N\mid n<x\}=n+1.
      \]\par
    \item
      \[
        (m\in\N)\land(n\in\N)\land(n<m)\Rightarrow(n+1\le m).
      \]\par
    \item Число $(n+1)\in\N$ непосредственно следует в $\N$ за натуральным
      числом $n$, т.~е. нет натуральных чисел $x$, удовлетворяющих условию
      $n<x<n+1$, если $n\in\N$.\par
    \item Если $n\in\N$ и $n\ne1$, то число $(n-1)\in\N$ и $(n-1)$
      непосредственно предшествует числу $n$ в $\N$, т.~е. нет натуральных
      чисел $x$, удовлетворяющих условию $n-1<x<n$, если $n\in\N$.\par
    \item В любом непустом подмножестве множества натуральных чисел имеется
      минимальный элемент.
  \end{enumerate}
\end{properties}

\subsection{Рациональные и иррациональные числа}

\begin{definition}{Целые числа}{integers}
  Объединение множества натуральных чисел, множества чисел, противоположных
  натуральным числам, и нуля называется множеством целых чисел и
  обозначается символом $\Z$.
\end{definition}

\begin{properties}{Арифметические свойства целых чисел}{integer-arithmetic}
  Сложение и умножение целых чисел не выводят за пределы множества $\Z$.
  Таким образом, $\Z$ есть абелева группа по отношению к операции сложения.
  По отношению к операции умножения множество $\Z$ и даже
  $\Z\setminus\{0\}$ не является группой, поскольку числа, обратные целым,
  не лежат в $\Z$ (кроме числа, обратного единице и минус единице).
\end{properties}

\begin{definition}{Делимость целых чисел}{integer-divisibility}
  В том случае, когда для чисел $m,n\in\Z$ число $k=m\cdot n^{-1}\in\Z$,
  т.~е. когда $m=k\cdot n$, где $k\in\Z$, говорят, что целое число $m$
  делится на целое число $n$ или кратно $n$, или что $n$ есть делитель $m$.
\end{definition}

\begin{definition}{Простое число}{prime-number}
  Число $p\in\N$, $p\ne1$, называется простым, если в $\N$ у него нет
  делителей, отличных от $1$ и $p$.
\end{definition}

\begin{theorem}{Основная теорема арифметики}{fundamental-arithmetic}
  Каждое натуральное число допускает и притом единственное (с точностью до
  порядка сомножителей) представление в виде произведения
  \[
    n=p_1\ldots p_k,
  \]
  где $p_1,\ldots,p_k$ — простые числа.
\end{theorem}

\begin{definition}{Взаимно простые числа}{coprime-integers}
  Числа $m,n\in\Z$ называются взаимно простыми, если у них нет общих
  делителей, отличных от $1$ и $-1$.
\end{definition}

\begin{corollary}{О делимости произведения взаимно простых чисел}{prime-divides-product}
  Если произведение $m\cdot n$ взаимно простых чисел $m,n$ делится на
  простое число $p$, то одно из чисел $m,n$ также делится на $p$.
\end{corollary}

\begin{definition}{Рациональные числа}{rational-numbers}
  Числа вида $m\cdot n^{-1}$, где $m,n\in\Z$, называются рациональными.
  Множество рациональных чисел обозначается знаком $\Q$.

  Число $q=m\cdot n^{-1}$ записывают также в виде отношения $m$ и $n$ или
  рациональной дроби $\frac{m}{n}$.
\end{definition}

\begin{properties}{Представление рационального числа}{rational-representation}
  Две упорядоченные пары $(m_1,n_1)$, $(m_2,n_2)$ задают одно и то же
  рациональное число тогда и только тогда, когда они пропорциональны,
  т.~е. существует число $k\in\Z$ такое, что, например,
  \[
    m_2=km_1,\qquad n_2=kn_1.
  \]
\end{properties}

\begin{definition}{Иррациональные числа}{irrational-numbers}
  Действительные числа, не являющиеся рациональными, называются
  иррациональными.
\end{definition}

\begin{proposition}{Существование иррационального числа}{sqrt-two-irrational}
  Существует положительное действительное число $s\in\R$, квадрат которого
  равен двум, и $s\notin\Q$.
\end{proposition}

\begin{definition}{Алгебраические и трансцендентные числа}{algebraic-transcendental}
  Вещественное число называется алгебраическим, если оно является корнем
  некоторого алгебраического уравнения
  \[
    a_0x^n+\ldots+a_{n-1}x+a_n=0
  \]
  с рациональными (или, что эквивалентно, с целыми) коэффициентами.

  В противном случае число называется трансцендентным.
\end{definition}

\subsection{Принцип Архимеда}

\begin{properties}{Свойство ограниченного подмножества $\N$}{bounded-natural-subset}
  В любом непустом ограниченном сверху подмножестве множества натуральных
  чисел имеется максимальный элемент.
\end{properties}

\begin{corollary}{О множествах натуральных и целых чисел}{natural-integer-bounds}
  \begin{enumerate}[label=\arabic*${}^{\circ}$.,start=2]
    \item Множество натуральных чисел не ограничено сверху.\par
    \item В любом непустом ограниченном сверху подмножестве множества целых
      чисел имеется максимальный элемент.\par
    \item В любом непустом ограниченном снизу подмножестве множества целых
      чисел имеется минимальный элемент.\par
    \item Множество целых чисел не ограничено ни сверху, ни снизу.
  \end{enumerate}
\end{corollary}

\begin{properties}{Принцип Архимеда}{archimedes-principle}
  Если фиксировать произвольное положительное число $h$, то для любого
  действительного числа $x$ найдется и притом единственное целое число $k$
  такое, что
  \[
    (k-1)h\le x<kh.
  \]
\end{properties}

\begin{corollary}{Следствия принципа Архимеда}{archimedes-corollaries}
  \begin{enumerate}[label=\arabic*${}^{\circ}$.,start=7]
    \item Для любого положительного числа $\eps$ существует натуральное
      число $n$ такое, что
      \[
        0<\frac1n<\eps.
      \]\par
    \item Если число $x\in\R$ таково, что $0\le x$ и для любого $n\in\N$
      выполнено $x<\frac1n$, то $x=0$.\par
    \item Для любых чисел $a,b\in\R$ таких, что $a<b$, найдется рациональное
      число $r\in\Q$ такое, что
      \[
        a<r<b.
      \]\par
    \item Для любого числа $x\in\R$ существует и притом единственное целое
      число $k\in\Z$ такое, что
      \[
        k\le x<k+1.
      \]
  \end{enumerate}
\end{corollary}

\begin{definition}{Целая и дробная части числа}{integer-fractional-parts}
  Указанное число $k$ обозначается $[x]$ и называется целой частью числа
  $x$. Величина
  \[
    \{x\}:=x-[x]
  \]
  называется дробной частью числа $x$. Итак, $x=[x]+\{x\}$, причем
  $\{x\}\ge0$.
\end{definition}

\subsection[Геометрическая интерпретация и вычислительные аспекты]{Геометрическая интерпретация множества действительных чисел и вычислительные аспекты операций с действительными числами}

\begin{definition}{Координата и числовая ось}{coordinate-axis}
  Число $f(x)$, соответствующее точке $x\in L$, называется координатой
  точки $x$. Прямую $L$ при наличии указанного соответствия
  $f\colon L\to\R$ называют координатной осью, числовой осью или числовой
  прямой. Ввиду биективности $f$ само множество $\R$ вещественных чисел
  также часто называют числовой прямой, а его элементы — точками числовой
  прямой.
\end{definition}

\clearpage
\begin{definition}{Числовые промежутки}{number-intervals}
  \[
    \begin{aligned}
      ]a,b[&:=\{x\in\R\mid a<x<b\} &&\text{— интервал }ab,\\
      [a,b]&:=\{x\in\R\mid a\le x\le b\} &&\text{— отрезок }ab,\\
      ]a,b]&:=\{x\in\R\mid a<x\le b\} &&\text{— полуинтервал }ab,
        \text{ содержащий конец }b,\\
      [a,b[&:=\{x\in\R\mid a\le x<b\} &&\text{— полуинтервал }ab,
        \text{ содержащий конец }a.
    \end{aligned}
  \]
  Интервалы, отрезки и полуинтервалы называются числовыми промежутками
  или просто промежутками. Числа, определяющие промежуток, называются его
  концами.

  Величина $b-a$ называется длиной промежутка $ab$. Если $I$ — некоторый
  промежуток, то длину его обозначают через $|I|$.

  Множества
  \[
    \begin{aligned}
      ]a,+\infty[&:=\{x\in\R\mid a<x\},&
      ]-\infty,b[&:=\{x\in\R\mid x<b\},\\
      [a,+\infty[&:=\{x\in\R\mid a\le x\},&
      ]-\infty,b]&:=\{x\in\R\mid x\le b\},
    \end{aligned}
  \]
  а также $]-\infty,+\infty[:=\R$, принято называть неограниченными
  промежутками.
\end{definition}

\begin{definition}{Окрестность точки}{neighborhood}
  Интервал, содержащий точку $x\in\R$, называется окрестностью этой точки.
  В частности, при $\delta>0$ интервал $]x-\delta,x+\delta[$ называется
  $\delta$-окрестностью точки $x$. Его длина $2\delta$.
\end{definition}

\begin{definition}{Модуль числа}{absolute-value}
  Функция
  \[
    |x|=
    \begin{cases}
      x,  & \text{при }x>0,\\
      0,  & \text{при }x=0,\\
      -x, & \text{при }x<0
    \end{cases}
  \]
  называется модулем или абсолютной величиной числа.
\end{definition}

\begin{definition}{Расстояние между числами}{distance-real}
  Расстоянием между $x,y\in\R$ называется величина $|x-y|$.
\end{definition}

\begin{properties}{Свойства расстояния и абсолютной величины}{absolute-value-properties}
  Расстояние неотрицательно и равно нулю только при совпадении $x$ и $y$;
  расстояние от $x$ до $y$ и от $y$ до $x$ одно и то же, ибо
  $|x-y|=|y-x|$; наконец, если $z\in\R$, то
  \[
    |x-y|\le|x-z|+|z-y|.
  \]

  Для любых чисел $x,y$ справедливо неравенство
  \[
    |x+y|\le|x|+|y|,
  \]
  причем равенство в нем имеет место в том и только в том случае, когда оба
  числа $x,y$ неотрицательны или неположительны.

  Кроме того,
  \[
    |x_1+\ldots+x_n|\le|x_1|+\ldots+|x_n|,
  \]
  причем равенство имеет место, если и только если все числа
  $x_1,\ldots,x_n$ одновременно неотрицательны или одновременно
  неположительны.
\end{properties}

\begin{definition}{Середина промежутка}{interval-center}
  Число $\frac{a+b}{2}$ называется серединой или центром промежутка с
  концами $a,b$, поскольку оно равноудалено от концов промежутка.
\end{definition}

\begin{definition}{Абсолютная и относительная погрешности}{errors}
  Если $x$ — точное значение некоторой величины, а $\widetilde{x}$ — известное
  приближенное значение той же величины, то числа
  \[
    \Delta(\widetilde{x}):=|x-\widetilde{x}|,
    \qquad
    \delta(\widetilde{x}):=
      \frac{\Delta(\widetilde{x})}{|\widetilde{x}|}
  \]
  называются соответственно абсолютной и относительной погрешностью
  приближения $\widetilde{x}$. Относительная погрешность при
  $\widetilde{x}=0$ не определена.
\end{definition}

\begin{proposition}{Погрешности арифметических операций}{arithmetic-errors}
  Если
  \[
    |x-\widetilde{x}|=\Delta(\widetilde{x}),
    \qquad
    |y-\widetilde{y}|=\Delta(\widetilde{y}),
  \]
  то
  \begin{align}
    \Delta(\widetilde{x}+\widetilde{y})
      &:=|(x+y)-(\widetilde{x}+\widetilde{y})|
        \le\Delta(\widetilde{x})+\Delta(\widetilde{y}),\label{eq:error-sum}\\
    \Delta(\widetilde{x}\widetilde{y})
      &:=|xy-\widetilde{x}\widetilde{y}|
        \le|\widetilde{x}|\Delta(\widetilde{y})
          +|\widetilde{y}|\Delta(\widetilde{x})
          +\Delta(\widetilde{x})\Delta(\widetilde{y}).\label{eq:error-product}
  \end{align}
  Если, кроме того,
  \[
    y\ne0,\qquad \widetilde{y}\ne0,
    \qquad
    \delta(\widetilde{y})=
      \frac{\Delta(\widetilde{y})}{|\widetilde{y}|}<1,
  \]
  то
  \begin{equation}
    \Delta\!\left(\frac{\widetilde{x}}{\widetilde{y}}\right)
      :=\left|\frac{x}{y}-\frac{\widetilde{x}}{\widetilde{y}}\right|
      \le
      \frac{|\widetilde{x}|\Delta(\widetilde{y})
        +|\widetilde{y}|\Delta(\widetilde{x})}{\widetilde{y}^{,2}}
      \cdot\frac{1}{1-\delta(\widetilde{y})}.
    \label{eq:error-quotient}
  \end{equation}
\end{proposition}

\begin{properties}{Оценки относительных погрешностей}{relative-errors}
  \begin{align}
    \delta(\widetilde{x}+\widetilde{y})
      &\le\frac{\Delta(\widetilde{x})+\Delta(\widetilde{y})}
        {|\widetilde{x}+\widetilde{y}|},\label{eq:relative-sum}\\
    \delta(\widetilde{x}\widetilde{y})
      &\le\delta(\widetilde{x})+\delta(\widetilde{y})
        +\delta(\widetilde{x})\delta(\widetilde{y}),\label{eq:relative-product}\\
    \delta\!\left(\frac{\widetilde{x}}{\widetilde{y}}\right)
      &\le\frac{\delta(\widetilde{x})+\delta(\widetilde{y})}
        {1-\delta(\widetilde{y})}.
    \label{eq:relative-quotient}
  \end{align}
\end{properties}

\begin{lemma}{О порядке числа по основанию $q$}{q-order-lemma}
  Если фиксировать число $q>1$, то для любого положительного числа
  $x\in\R$ найдется и притом единственное целое число $k\in\Z$ такое, что
  \[
    q^{k-1}\le x<q^k.
  \]
\end{lemma}

\begin{definition}{Порядок числа по основанию $q$}{q-order}
  Пусть $q>1$ и $x>0$. Число $p$, удовлетворяющее соотношению
  \[
    q^p\le x<q^{p+1},
  \]
  называется порядком числа $x$ по основанию $q$ или (при фиксированном
  $q$) просто порядком числа $x$.
\end{definition}

\begin{definition}{$q$-ичная позиционная запись}{q-ary-representation}
  Каждому положительному числу $x\in\R$ взаимно однозначно сопоставляется
  символ вида $\alpha_p\ldots\alpha_0,\ldots$, если $p\ge0$, или
  $0,0\ldots0\alpha_p\ldots$, если $p<0$, причем
  $\alpha_k\in\{0,1,\ldots,q-1\}$, как угодно далеко встречаются знаки,
  отличные от $q-1$, и
  \[
    r_n=\alpha_pq^p+\ldots+\alpha_{p-n}q^{p-n},
    \qquad
    r_n\le x<r_n+\frac{1}{q^{n-p}}.
  \]
  Он называется $q$-ичной позиционной записью числа $x$; цифры, входящие в
  символ, называют знаками; позиции знаков относительно запятой называются
  разрядами.
\end{definition}

\begin{definition}{Сечение множества рациональных чисел}{dedekind-cut}
  Сечением в множестве $\Q$ рациональных чисел называется разбиение $\Q$
  на два не имеющих общих элементов множества $A$, $B$ таких, что
  \[
    \forall a\in A\;\forall b\in B\;(a<b).
  \]
\end{definition}

% -----------------------------------------------------------------------------
\clearpage
\lecture[Основные леммы полноты $\R$]{Основные леммы, связанные с полнотой множества действительных чисел}

\subsection{Лемма о вложенных отрезках (принцип Коши—Кантора)}

\begin{definition}{Последовательность}{sequence}
  Функцию $f\colon\N\to X$ натурального аргумента называют
  последовательностью или, полнее, последовательностью элементов множества
  $X$.

  Значение $f(n)$ функции $f$, соответствующее числу $n\in\N$, часто
  обозначают через $x_n$ и называют $n$-м членом последовательности.
\end{definition}

\begin{definition}{Последовательность вложенных множеств}{nested-sets}
  Пусть $X_1,X_2,\ldots,X_n,\ldots$ — последовательность каких-то множеств.
  Если
  \[
    X_1\supset X_2\supset\ldots\supset X_n\supset\ldots,
  \]
  т.~е. $\forall n\in\N\;(X_n\supset X_{n+1})$, то говорят, что имеется
  последовательность вложенных множеств.
\end{definition}

\begin{lemma}{Коши—Кантора}{nested-intervals}
  Для любой последовательности
  \[
    I_1\supset I_2\supset\ldots\supset I_n\supset\ldots
  \]
  вложенных отрезков найдется точка $c\in\R$, принадлежащая всем этим
  отрезкам.

  Если, кроме того, известно, что для любого $\eps>0$ в последовательности
  можно найти отрезок $I_k$, длина которого $|I_k|<\eps$, то $c$ —
  единственная общая точка всех отрезков.
\end{lemma}

\subsection{Лемма о конечном покрытии (принцип Бореля—Лебега)}

\begin{definition}{Покрытие множества}{cover}
  Говорят, что система $S=\{X\}$ множеств $X$ покрывает множество $Y$, если
  \[
    Y\subset\bigcup_{X\in S}X,
  \]
  т.~е. если любой элемент $y$ множества $Y$ содержится по крайней мере в
  одном из множеств $X$ системы $S$.

  Подмножество множества $S=\{X\}$, являющегося системой множеств, называют
  подсистемой системы $S$.
\end{definition}

\begin{lemma}{Бореля—Лебега}{finite-cover}
  В любой системе интервалов, покрывающей отрезок, имеется конечная
  подсистема, покрывающая этот отрезок.
\end{lemma}

\subsection{Лемма о предельной точке (принцип Больцано—Вейерштрасса)}

\begin{definition}{Предельная точка множества}{limit-point}
  Точка $p\in\R$ называется предельной точкой множества $X\subset\R$, если
  любая окрестность этой точки содержит бесконечное подмножество множества
  $X$.
\end{definition}

\begin{properties}{Критерий предельной точки}{limit-point-criterion}
  Это условие равносильно тому, что в любой окрестности точки $p$ есть по
  крайней мере одна не совпадающая с $p$ точка множества $X$.
\end{properties}

\begin{lemma}{Больцано—Вейерштрасса}{bolzano-weierstrass}
  Всякое бесконечное ограниченное числовое множество имеет по крайней мере
  одну предельную точку.
\end{lemma}

% -----------------------------------------------------------------------------
\clearpage
\lecture{Счётные и несчётные множества}

\subsection{Счётные множества}

\begin{definition}{Счётное множество}{countable-set}
  Множество $X$ называется счётным, если оно равномощно множеству $\N$
  натуральных чисел, т.~е.
  \[
    \card X=\card\N.
  \]
\end{definition}

\begin{proposition}{О счётных множествах}{countable-sets-proposition}
  \begin{enumerate}[label=\alph*)]
    \item Бесконечное подмножество счётного множества счётно.\par
    \item Объединение множеств конечной или счётной системы счётных множеств
      есть множество счётное.
  \end{enumerate}
\end{proposition}

\begin{definition}{Не более чем счётное множество}{at-most-countable}
  Если про множество известно, что оно либо конечно, либо счётно, то говорят,
  что оно не более чем счётно. Равносильная запись:
  \[
    \card X\le\card\N.
  \]
\end{definition}

\begin{corollary}{Следствия утверждения о счётных множествах}{countable-corollaries}
  \begin{enumerate}[label=\arabic*)]
    \item $\card\Z=\card\N$.\par
    \item $\card\N^2=\card\N$.\par
    \item $\card\Q=\card\N$, т.~е. множество рациональных чисел счётно.\par
    \item Множество алгебраических чисел счётно.
  \end{enumerate}
\end{corollary}

\subsection{Мощность континуума}

\begin{definition}{Числовой континуум}{continuum}
  Множество $\R$ действительных чисел называют также числовым континуумом,
  а его мощность — мощностью континуума.
\end{definition}

\begin{theorem}{Кантора}{cantor-real}
  \[
    \card\N<\card\R.
  \]
\end{theorem}

\begin{corollary}{Следствия теоремы Кантора}{cantor-real-corollaries}
  \begin{enumerate}[label=\arabic*)]
    \item $\Q\ne\R$ и существуют иррациональные числа.\par
    \item Существуют трансцендентные числа, поскольку множество
      алгебраических чисел счётно.
  \end{enumerate}
\end{corollary}
